\section{Round-ten reconstruction: closed nullspaces, common transfer bundles, and typed contractions}
\label{sec:r10-c2}

Every platform retains its label.  Cross-platform statements require an
explicit scaling map.  Within one platform, path potentials live in a weighted
dynamical algebra \(\mathscr A_W\) invariant under time shifts.

\subsection{Two genuine closed linear nullspaces}

Define
\[
 \mathscr N_{\rm int}=
 \overline{\operatorname{span}}
 \{U-U\circ\Theta_t:U\in\mathscr A_W,
                    t\in\mathbb Q_+\},
\]
\[
 \mathscr N_{\rm pr}=
 \overline{\operatorname{span}}
 \bigl(\mathbb R\mathbf1+\mathscr N_{\rm int}\bigr).
\]
The closures are in the weighted strict topology.

\begin{lemma}[Platform-wise uniform Ces\`aro tightness]
\label{lem:r10-c2-cesaro}
On each platform, the chosen weight satisfies
\[
 \sup_{T\ge1}\frac1T\int_0^T
 \int W\circ\Theta_t\,d|\mu|\,dt
 \le C\int(1+W)d|\mu|
\]
for every weighted finite signed measure \(\mu\).  Consequently the Ces\`aro
averages of \((\Theta_t)_\#\mu\) are uniformly weighted-tight.
\end{lemma}

\begin{proof}
For the compact Sinai platform, the path weight is controlled by the
exponential return/roof Lyapunov function, whose expected drift is negative
outside a compact set.  For hard spheres, mass and energy are conserved and
the contact-count weight has a uniform short-time exponential moment from B2.
Integrating the respective Foster--Lyapunov inequalities yields the displayed
uniform averaged bound.  Weighted Prokhorov compactness gives tightness.
\end{proof}

\begin{theorem}[Invariant separator and Hausdorff cotangents]
\label{thm:r10-c2-null}
A potential has zero integral under every invariant weighted probability if
and only if it belongs to \(\mathscr N_{\rm int}\).  Two potentials have the
same long-time pressure if and only if their difference belongs to
\(\mathscr N_{\rm pr}\).  Both quotients are Hausdorff, and their continuous
duals are countably additive invariant Radon measures with the appropriate
normalization.
\end{theorem}

\begin{proof}
Coboundaries integrate to zero and constants shift pressure.  Conversely,
separate a class outside the closed linear span by a weighted Radon functional
using the strict dual theorem.  Average that functional over rational times.
Lemma~\ref{lem:r10-c2-cesaro} gives a weighted-tight subsequence; its limit is
invariant and preserves the separating pairing because the averaged
coboundary error is an endpoint term divided by time.  Jordan decomposition
and normalization give invariant probabilities.  The pressure statement
follows from the variational principle and the constant direction.
\end{proof}

\subsection{A common transfer Banach bundle}

For a regular phase chart \(\eta\mapsto\mathcal L_\eta\), all transfer
operators act on the common strong/weak Banach bundle furnished by A1 or A2,
not on mutually identified \(L^2(\nu_\eta)\) spaces.  Let \(\Pi_\eta\) be the
simple Riesz projector and \(S_\eta\) its reduced resolvent.

\begin{theorem}[Kato transport without infinite-path Radon--Nikodym densities]
\label{thm:r10-c2-kato}
The differential equation
\[
 U_\eta'=[\Pi_\eta',\Pi_\eta]U_\eta,
 \qquad U_0=I,
\]
defines bounded Kato parallel transport on the common transfer bundle and
satisfies \(U_\eta\Pi_0=\Pi_\eta U_\eta\).  Covariant derivatives of the
leading eigenfunction, eigenfunctional, pressure, spectral projector and any
finite resolved subspace are finite sums of \(\mathcal L_\eta'\),
\(S_\eta\), and their iterates.  No Radon--Nikodym derivative between distinct
infinite-volume ergodic path measures is used.
\end{theorem}

\begin{proof}
The Riesz formula on the common Banach scale gives a continuously
differentiable finite-rank projector.  Kato's commutator equation has a unique
bounded evolution and differentiating
\(U_\eta^{-1}\Pi_\eta U_\eta\) shows it is constant.  Differentiating the
eigen-equations and applying the reduced resolvent gives the response words.
All objects remain in the common transfer space even when the corresponding
stationary path measures are mutually singular.
\end{proof}

\subsection{Resolved memory on the transported subspace}

Let \(P_\eta=U_\eta P_0U_\eta^{-1}\) be the transported finite resolved
projection in the common Hilbert realization of the A4 Doob phase.  Put
\[
 C_\eta(z)=P_\eta(z-L_\eta)^{-1}P_\eta,
 \qquad
 \widehat K_\eta(z)=zP_\eta-P_\eta L_\eta P_\eta-C_\eta(z)^{-1}.
\]

\begin{theorem}[Covariant compressed resolvent and memory]
\label{thm:r10-c2-memory}
The covariant derivative of \(C_\eta\) is the resolvent Duhamel insertion plus
the two transported projection derivatives.  Differentiating the displayed
identity gives the covariant memory response, including derivatives of the
phase generator and resolved subspace.  It agrees with first differentiating
the A4 history pressure and then compressing.
\end{theorem}

\begin{proof}
Differentiate
\(P_\eta(z-L_\eta)^{-1}P_\eta\) in the common bundle and use
\(D(z-L)^{-1}=(z-L)^{-1}(DL)(z-L)^{-1}\).  Matrix inversion gives
\(D(C^{-1})=-C^{-1}(DC)C^{-1}\).  The A4 pressure tangent is the same
propagated correlation operator before Laplace transformation, so Fubini and
normal resolvent bounds prove the commutation.  The sign is the causal
Volterra sign of A4.
\end{proof}

\subsection{Optional projections and nonlinear contractions}

\begin{theorem}[Stable optional-projection limit]
\label{thm:r10-c2-optional}
Finite deterministic likelihood martingales projected onto the A4 past
filtration converge stably, jointly with the enhanced path tangent, to the
optional projection of the limiting diffusion likelihood.  The result holds
for every bounded terminal cylinder and every regular phase chart.
\end{theorem}

\begin{proof}
A4 gives locally uniform convergence of conditional future kernels on compact
past histories and uniform exponential moments.  Apply the kernels to the
finite terminal likelihoods and use stable convergence of the enhanced path.
Uniform integrability permits conditional expectations to pass to the limit.
The limit is the diffusion optional projection by uniqueness of regular
conditional expectation.
\end{proof}

\begin{theorem}[Typed nonlinear contraction and tangent chain rule]
\label{thm:r10-c2-main}
Let \(\mathcal C:X\to Y\) be a continuous nonlinear observable map with
Fr\'echet derivative on the regular rate sublevel.  The contracted pressure is
\[
 Q_{\mathcal C}(\varphi)=Q(\varphi\circ\mathcal C),
\]
not \(Q(\mathcal C^*\varphi)\) for a fictitious nonlinear adjoint.  At an
exposed phase \(x_*\),
\[
 DQ_{\mathcal C}(\varphi)h
 =DQ(\varphi\circ\mathcal C)
   [D\varphi(\mathcal Cx_*)D\mathcal C(x_*)h],
\]
and the Hessian contains the covariance pullback plus the second derivatives
of \(\varphi\) and \(\mathcal C\).  The rate is the usual nonlinear
contraction \(I_{\mathcal C}(y)=\inf_{\mathcal Cx=y}I(x)\).
\end{theorem}

\begin{proof}
The finite identity
\(e^{\mu\varphi(\mathcal C(X_\varepsilon))}\) is exactly the original source
functional composed with \(\mathcal C\).  The contraction principle gives the
rate.  Fr\'echet chain rules on the regular chart give the first and second
responses, while the covariance term is supplied by A4/B3.  This formulation
is typed for a nonlinear map and requires no nonlinear linear adjoint.
\end{proof}
